\section{Limitations and Roadmap}\label{sec:limitations}

The present report's contributions are bounded in five respects.
We state the boundaries and indicate the remediation path for
each.

\paragraph{INLINE write path.}
The INLINE scheme implements only the read path. The write side
returns \texttt{-EOPNOTSUPP} at \texttt{write\_begin},
\texttt{write\_end}, and \texttt{writepages}. Adding the write
path requires implementing per-block RS encoding on the writeback
path, RS journal entry on successful write, and proper
interaction with the bitmap allocator and the inode metadata
update. The estimated incremental kernel module size is
\(\sim\)250 lines of C, comfortably within the remaining 2{,}802
LOC of the 5{,}000-line audit budget. INLINE write is part of the
v2.x roadmap, post the present technical report and Zenodo
upload.

\paragraph{Saturation falsification (Theorem v2.2).}
The empirical section corroborates \cref{thm:soundness-v21} on a
single-symbol-error event well within the
correction capacity. \cref{thm:saturation-v22} (saturation
observability) is stated in the present report but not yet
exercised by an active falsification protocol. A saturation
protocol would force \(>8\) symbol errors on a single subblock,
verify that \(\mathcal{G}\) returns \(\bot\), and verify that the
journal contains a \texttt{UNCORRECTABLE}-flagged entry with the
correct \((\text{block}, \text{subblock})\) coordinates. The
protocol is straightforward to derive from the Stage 3
\texttt{dd}-based corruption template, by injecting nine or more
distinct symbol-position flips in one subblock; we mark this as
the immediate post-Zenodo experimental task.

\paragraph{Companion fault injection diagnostic.}
The companion fault injection module RadFI v0.1.2 exhibited a
diagnostic anomaly during the present session: the kprobe
pre-handler did not fire despite the kprobe being armed at the
exact symbol address and the symbol being demonstrably called by
the workload (verified by native kprobe events recorded directly
through the tracefs interface). The hypotheses ruled out by
direct inspection include kprobe blacklist, function inlining,
\texttt{submit\_bh} versus \texttt{submit\_bio\_noacct} dispatch
mismatch, and BEAMFS hook flag set. The remaining hypotheses
under investigation include
\texttt{CONFIG\_KPROBES\_ON\_FTRACE} dispatch interaction and
\texttt{bd\_dev} encoding in the targeting algebra. Diagnosis is
deferred to the post-Zenodo continuation; the empirical results
of the present report use \texttt{dd}-based corruption as a
direct, reproducible substitute. RadFI's own methodology paper
remains the formal reference~\cite{desbrieres2026radfi}.

\paragraph{Verified compilation.}
The present report's recovery operator is defined and proved at
the mathematical level (\cref{def:G,thm:soundness-v21,thm:saturation-v22});
the kernel implementation is paired with the operator definition
through the traceability table of \cref{tab:m2c-trace}. We do not
claim that the implementation satisfies the operator definition
in any machine-checked sense. Closing this gap requires either an
extraction-based approach in the FSCQ
tradition~\cite{chen2015fscq} or a refinement-based approach in
the seL4 tradition~\cite{klein2009sel4}, applied to the BEAMFS
recovery operator and the \texttt{lib/reed\_solomon} kernel
substrate~\cite{kernelrslib}. We mark this as v3+ work.

\paragraph{Linux fs-devel submission.}
This report is intended for Zenodo upload and citation in
academic and industrial contexts where electromagnetic resilience
is a stated design constraint. The Linux kernel-tree submission
to \texttt{linux-fsdevel} is a separate workstream with its own
prerequisites: completion of the INLINE write path, an
\texttt{xfstests} subset passing, strict \texttt{checkpatch.pl}
on every modified file, formal LOC audit against the certification
regimes (DO-178C, ECSS-E-ST-40C, IEC\,61508) targeting the
sub-5{,}000-LOC budget, and a cover letter aligned with the
threading model preferred by Linux maintainers. The estimated
calendar time is two to three months from the date of the present
report, conditional on review cycle responsiveness.
